% 04_methods -- front lane.
% src: juliet_full/static/README.md; juliet_full/llm/README.md; juliet_full/llm/prompts/; juliet_full/llm/l4/leak.dl;
%      juliet_full/gpu/README.md; juliet_full/decision_ft/README.md; juliet_full/analysis/score.py; juliet_full/STATUS.md
\section{Approaches and Setup}
\label{sec:methods}

Table~\ref{tab:methods} lists every detector.
All of them see the same \Deb{} unit text and answer one question: does the executed unit leak heap memory?
Everything that could be tuned (checker mapping, prompts, examples, thresholds, training checkpoints) was fixed on dev before any eval run.

% front lane. src: juliet_full/STATUS.md (versions), juliet_full/lane_reports/2026-11-23_static_stage1.md (checker mapping),
% juliet_full/llm/README.md, juliet_full/gpu/README.md, juliet_full/decision_ft/README.md
\begin{table}[t]
\centering
\caption{Detectors. Strategies: L1 direct, L2 few-shot, L3 with Clang SA warnings, L4 LLM facts + Datalog. Dynamic tools execute the code and serve only as a reference.}
\label{tab:methods}
\footnotesize
\setlength{\tabcolsep}{3pt}
\begin{tabularx}{\columnwidth}{@{}l>{\raggedright\arraybackslash}X l@{}}
\toprule
Detector & Technique / configuration & Runs \\
\midrule
grep rule & no \code{free}/\code{delete} $\Rightarrow$ leak & -- \\
Cppcheck 2.22.0 & patterns + value flow; \code{memleak*}, \code{leak*} checks & -- \\
Clang SA 23.1.6 & path-sensitive symbolic execution; \code{unix.Malloc}, \code{cplusplus.NewDeleteLeaks} & -- \\
Infer 1.3.0 & Pulse (incorrectness separation logic); \code{MEMORY\_LEAK\_C/CPP} & -- \\
CodeQL 2.27.5 & Datalog dataflow queries; never-freed, may-not-be-freed, realloc & -- \\
LineVul & CodeBERT fine-tuned on Big-Vul; zero-shot, 512 tokens & -- \\
\midrule
DeepSeek-V4.1-Flash & API LLM, temperature 0, no thinking & L1--L4 \\
Qwen3.5-4B & local LLM (vLLM), temperature 0, no thinking & L1--L4 \\
\midrule
Jev 1.14.0 & hosted decision model, one \code{noul} question & L1--L3 \\
Intern-Decision-4B & local decision model, placeholder logits & L1--L3 \\
LoRA decision (ours) & Qwen3.5-4B + LoRA, binary head, 3-fold CV & L1, L3 \\
\midrule
ASan/LSan & GCC 13.3, \code{use\_stacks=0} & ref. \\
Valgrind 3.28.0 & Memcheck, definite + indirect & ref. \\
dynamic-enhanced & ASan/LSan + globals + \code{realloc} fault injection, 24 runs & ref. \\
\bottomrule
\end{tabularx}
\end{table}


\subsection{Static analyzers and the dynamic reference}
\label{sec:m-static}
Each analyzer runs on the source files of a unit.
A unit is a leak iff the analyzer reports at least one mapped finding in the unit's own files.
We chose the mapping on the dev units only; for every tool the union of its leak checkers had the best dev macro-F1 among rule subsets. % src: juliet_full/static/README.md (verdict mapping decisions)
Table~\ref{tab:methods} lists the mapped checkers.
Two settings mattered.
Clang needs \code{-analyzer-opt-analyze-headers}: without it, the path-sensitive engine skipped every function in an included source file, and dev recall on flows 51--84 was 0.
Infer gets each source file as its own compilation entry for the same reason. % src: static/README.md (dev_v1; per_file 0.710 -> 0.784)
CodeQL builds one database per batch of 50 units.
Each unit's \code{sym\_N} names are renamed per unit at compile time, and no batch holds both units of a group.
A tool error or timeout is an invalid prediction.

ASan/LSan and Valgrind use the settings of \S\ref{sec:verify}.
The \emph{dynamic-enhanced} reference covers the cases plain LSan misses by running every unit 24 times: three configurations (plain; \code{use\_globals=0}, so memory held only by a global at exit counts; the first \code{realloc} returns \code{NULL}) times eight fixed values of \code{time()}, which cover all branch combinations of variant~12. % src: juliet_full/data/README.md (dynamic_enhanced)
A unit is a leak if any run reports one.
On all good units of \Deb, \Dead{} and \Plac{} it raises no alarm.

\subsection{Prompting strategies L1--L4}
\label{sec:m-strategies}
DeepSeek and Qwen receive identical prompts (Appendix~\ref{app:prompts}).
The system prompt states the task and our leak definition from \S\ref{sec:background}.
It also describes the Juliet support code (print helpers, \code{ALLOCA}).
L1--L3 ask for one compact JSON object, \code{\{"label": ..., "confidence": ...\}}; the API runs use JSON mode.
\begin{itemize}
\item \textbf{L1} (direct) shows the unit in a code fence.
\item \textbf{L2} (few-shot) adds four dev units with their verdicts, two bad and two good, one per dev flow category, selected by a deterministic rule (shortest unit per slot). Both units of these four groups are excluded from every run.
\item \textbf{L3} (evidence) adds the Clang SA warnings of the frozen \code{full\_v2} run, with line numbers relative to the shown unit, or a sentence saying that there are none.
\item \textbf{L4} (neuro-symbolic) asks the model only for facts, not for a verdict. The schema has ten relations: \code{alloc}, \code{realloc}, \code{free}, \code{assign}, \code{pass}, \code{retval}, \code{ret}, \code{call}, \code{global} and \code{dead}. Every fact carries the line of its statement, and facts whose line does not contain the statement are dropped. A Datalog program for Souffl\'e~\cite{jordan2016souffle} (\code{leak.dl}, Appendix~\ref{app:datalog}) computes the functions reachable from \code{main} and a flow-insensitive, Andersen-style points-to relation. It then derives \code{leak(o, never\_freed)} for an allocated object that no live \code{free} reaches, and \code{leak(o, realloc\_fail)} for \code{p = realloc(p, ...)}. Facts on lines marked \code{dead} are ignored. The unit is a leak iff a leak is derived.
\end{itemize}
Prompt files are hashed, and an eval run refuses to start unless the hash equals the frozen value.
Decoding uses temperature 0 with thinking disabled; a run stops if a reasoning trace appears.
An output that cannot be parsed, an API error, or a timeout is counted as wrong, never retried with a different prompt. % src: juliet_full/analysis/score.py:9-10

\noindent\textbf{Models.}
DeepSeek-V4.1-Flash~\cite{deepseek2026v4} is called through its official API~\cite{deepseekdocs}.
Qwen3.5-4B~\cite{qwen35card} runs locally on vLLM~\cite{kwon2023vllm} (BF16, one RTX 3090).
The API budget was a hard cap of \yen10 for DeepSeek (the ledger stops at \yen9.5) and \$4 for Jev. % src: juliet_full/llm/README.md (Budget: DeepSeek CNY 9.5, Jev USD 4.0)
DeepSeek runs only in off-peak hours, which are half price.
Under this budget, DeepSeek covers all eval units on \Deb{} with L1--L4. On \Dead{} it covers \code{api\_subset} with L1--L3 and \code{nested25} with L4; on \Plac{} (L1, L3) and \Orig{} (L1) it covers \code{nested25}. % src: juliet_full/llm/results/deepseek_flash/*/ (line counts)

\subsection{Decision models}
\label{sec:m-decision}
\noindent\textbf{Jev} (hosted, version 1.14.0~\cite{typesafe2026jev}) gets the same L1--L3 content as its \emph{state} and one yes/no (\code{noul}) question with our leak criteria.
It returns a probability.
The primary threshold is 0.5; a secondary threshold per strategy, chosen on dev by macro-F1 (L1 0.58, L2 0.635, L3 0.585), is reported separately. % src: juliet_full/llm/results/jev_thresholds.json

\noindent\textbf{Intern-Decision-4B}~\cite{zhu2026interndecision} runs zero-shot on our GPU server with the same Jev prompts.
It reads the logits of one forward pass at a decision placeholder, as in its model card, and returns calibrated option probabilities.
We use threshold 0.5 and a secondary threshold frozen on 652 dev units before any eval call (L1 0.823, L2 0.673, L3 0.539). % src: juliet_full/gpu/README.md:139-140

\noindent\textbf{Our LoRA decision model} starts from Qwen3.5-4B (BF16, frozen) and adds LoRA~\cite{hu2022lora} adapters (rank 16, $\alpha=32$, dropout 0.05) to all linear projections of the language model: attention, linear attention and MLP. This gives 32,464,896 trainable parameters. % src: juliet_full/decision_ft/README.md §3
The head is a $2\times2560$ linear layer initialized from the output-embedding rows of the tokens \code{Yes} and \code{No}.
It reads the final hidden state at the last prompt position, and the score is $s=\mathrm{logit}_\mathrm{Yes}-\mathrm{logit}_\mathrm{No}$.
The input is the L1 or L3 prompt followed by ``Does the executed code unit leak heap memory? Answer Yes or No.''
The loss is class-balanced cross-entropy plus a group ranking term,
\[
\mathcal{L}=\mathrm{CE}+\lambda\cdot\frac{1}{|G|}\sum_{g\in G}-\log\sigma\!\left(\frac{s_{\mathrm{bad}(g)}-s_{\mathrm{good}(g)}}{\tau}\right),
\]
with $\tau=1$ and $\lambda=1$. Each micro-batch holds whole groups.
We train with AdamW (LoRA lr $10^{-4}$, head lr $10^{-3}$, 5\% warmup, cosine decay) for at most three epochs, at eight groups per step.
The model is evaluated on dev every half epoch.
We keep the checkpoint with the highest dev AUROC (ties: lower dev loss) and set the threshold on dev scores to maximize macro-F1.
The adapter is merged into the weights for prediction, and the reported numbers come from the merged model.
Training uses three-fold cross-validation over the eval variants (\S\ref{sec:split}).
Each fold model trains on two folds, uses the 330 dev groups only for selection, and predicts the held-out fold.
The six jobs (two inputs $\times$ three folds) took \res{x}{ft_hours_total} GPU hours.

\noindent\textbf{LineVul}~\cite{fu2022linevul} uses the public Big-Vul checkpoint zero-shot, with inputs truncated to 512 tokens and threshold 0.5.

\subsection{Runner, metrics and measurement}
\label{sec:m-metrics}
Every method writes one JSON line per unit in a common schema (unit, method, strategy, model, layer, prediction, score, latency, status), read by a single scorer.
Paid calls go to an append-only cost ledger with per-run caps, and DeepSeek runs pause during peak hours.

\noindent\textbf{Metrics.}
The positive class is \emph{leak}.
We report its precision, recall and F1, and macro-F1 over both classes as the headline, end to end with invalid outputs counted wrong.
For confidence intervals we use a percentile bootstrap that resamples the 42 families with replacement (2,000 iterations)~\cite{efron1979}, because units of one family share code.
Paired comparisons use McNemar's test~\cite{mcnemar1947} on per-unit correctness, exact binomial when there are at most 25 discordant pairs.
We apply a Holm correction~\cite{holm1979} across the planned comparisons.
For score-based methods we also report AUROC.
Robustness is the change from \Deb{} to \Dead{} and \Plac{} on the same units: \emph{lost} leaks (bad units detected on \Deb{} but missed on \Dead) and \emph{gained} false alarms (good units that were correct on \Deb{} but are flagged on \Dead).

\noindent\textbf{Latency and cost.}
Latency is wall-clock time per unit at concurrency~1 on the 300-unit \code{latency300} subset, with median and p90.
It includes the analyzers' own compilation; CodeQL is reported per unit (first 20 units) and amortized over batches of 50.
Latency runs hold a host-wide lock and start only when the host is quiet (1-minute load below 2.5).
Cost is the API charge from the ledger in CNY for DeepSeek and USD for Jev; local models are reported in GPU time.

\noindent\textbf{Hardware.}
Clients, static and dynamic tools run on a Dell OptiPlex 5000 (Intel i5-12500, 12 threads, 31\,GB RAM). % src: lscpu / free -g on dell
Qwen3.5-4B, Intern-Decision-4B and the LoRA model run on a server with two RTX 3090 GPUs.
Tool versions: Clang 23.1.6, Infer 1.3.0, CodeQL 2.27.5, Cppcheck 2.22.0, Valgrind 3.28.0, Souffl\'e 2.5, GCC 13.3. % src: juliet_full/STATUS.md:33; gcc --version
